\chapter*{Considerações finais}
\addcontentsline{toc}{chapter}{Considerações finais}
\label{cap:conclusao}

% =====================================================================
% PAPEL DESTAS CONSIDERAÇÕES: retomar a pergunta e a hipótese à luz dos
% resultados (Caps. 5 e 6), sintetizar contribuições, declarar limitações
% e desdobramentos. NÃO introduz dado novo: todo número aqui já apareceu
% em capítulo anterior, que é citado.
% =====================================================================

Esta dissertação partiu de uma pergunta prática, a de como apoiar o professor de Matemática do
Ensino Médio na geração de questões com modelos de linguagem garantindo correção matemática
verificável e adequação didática reconhecida. A resposta proposta foi um sistema multiagente em que
um modelo de linguagem redige a questão, um verificador simbólico recalcula o gabarito de forma
independente, um crítico didático avalia a questão segundo uma rubrica fundamentada e um orquestrador
decide entre aprovar, devolver para revisão ou descartar. O percurso de construção e avaliação desse
sistema permite agora retomar a pergunta com mais precisão do que no início.

O sistema mostrou-se capaz de produzir questões utilizáveis. No acervo de 90 especificações que
cobriu as quinze habilidades do recorte, 85 terminaram em questão aprovada, a maior parte delas na
primeira iteração, e o custo e o tempo medidos são compatíveis com o uso cotidiano de um professor
(Seção~\ref{sec:custo}). Quando essas questões chegaram a professores da rede pública, oito em cada
dez decisões as aceitaram com ou sem ajuste, e nenhuma recusa foi motivada por erro matemático
(Capítulo~\ref{cap:avaliacao}). Do ponto de vista do produto, portanto, o objetivo de oferecer ao
professor rascunhos de qualidade alinhados à BNCC foi alcançado.

A hipótese arquitetural, por sua vez, recebe uma resposta mais matizada, e é nela que se encontram os
resultados mais interessantes do trabalho. Com o modelo usado na geração do acervo, o ciclo de revisão
trabalhou sobretudo na dimensão didática, já que a grande maioria das revisões foi pedida pelo Crítico.
O Verificador, por outro lado, não encontrou nenhum gabarito matematicamente errado nas execuções
analisadas, e as suas reprovações ou eram indevidas ou apontavam inconsistências na própria
formalização. A sua contribuição apareceu de outra forma, como garantia declarada, uma vez que 72 das
85 questões chegaram ao professor com o gabarito recalculado de modo independente e com essa
informação registrada junto da questão. Ao mesmo tempo o único erro matemático do acervo atravessou o
sistema inteiro e um professor atento, porque estava na tradução do enunciado para a formalização e
não no cálculo (Seção~\ref{sub:rodada1-erros}). Esse episódio confirma a premissa do erro silencioso,
já que o gabarito errado passou por quem fez a conta certa na margem da folha, e ao mesmo tempo
delimita com exatidão o alcance da verificação simbólica. Ela elimina o cálculo errado sobre um
problema bem formalizado e não protege contra a leitura errada do problema.

Daí decorre a principal contribuição conceitual da dissertação, que é tratar a garantia como algo a
ser declarado com precisão e não como selo genérico de qualidade. O sistema distingue o que cada
habilidade admite de conferência automática do que cada questão de fato recebeu, e diz ao professor
em palavras simples o que o rótulo mais forte significa e o que ele não significa
(Seção~\ref{sec:garantia}). Em vez de prometer um sistema que não erra, o trabalho oferece um sistema
que torna visível onde o erro ainda pode estar, o que é coerente com a compreensão do professor como
mediador indispensável defendida na Seção~\ref{sec:papel-mediador}.

Outras contribuições merecem registro. A primeira é a especificação de entrada que parte da habilidade
da BNCC e não do tema, acompanhada das exigências que a questão precisa exibir para realizar a
habilidade, o que transforma o alinhamento curricular de rótulo em critério verificável pelo Crítico.
A segunda é o tipo de verificação por propriedades, que estende a conferência simbólica a habilidades
de generalização de padrões e de associação entre progressões e funções, nas quais não há um valor a
comparar. A terceira é de natureza metodológica. Os testes automatizados, por mais numerosos, não
encontraram os defeitos que importavam, e foram as execuções com provedor real que revelaram onze
deles no Verificador, todos do mesmo lado da balança, o de reprovar gabaritos corretos
(Seção~\ref{sec:reprodutibilidade}). Para quem desenvolve sistemas desse tipo, a lição é que um
verificador que reprova o que está certo não é inofensivo, porque induz o gerador a piorar uma questão
que estava correta.

O estudo tem limitações que precisam ser ditas com clareza. A mais importante é que o painel formal,
desenhado para testar a hipótese com itens comuns, desenho cego e medida de concordância, não foi
aplicado no prazo desta dissertação (Seção~\ref{sec:painel-formal}). A avaliação realizada é formativa,
com cinco professores que viram lotes diferentes, e descreve o que aconteceu sem estimar propriedades
do sistema. Além disso todas as gerações analisadas usaram um único modelo comercial, de modo que nada
se pode afirmar ainda sobre modelos menores ou abertos. A fidelidade da formalização ao enunciado não é
conferida por nenhum agente, e o sistema produz apenas texto, sem as figuras que três dos cinco
professores pediram. O recorte de quinze habilidades, embora coerente com uma unidade que a própria
BNCC reconhece, deixa de fora a geometria, a estatística e a probabilidade. Por fim, os professores
mostraram trabalhar com uma referência curricular diferente da usada no catálogo, ponto que precisa ser
esclarecido porque a habilidade declarada é a âncora do alinhamento.

Essas limitações apontam os desdobramentos mais imediatos. O primeiro é aplicar o painel formal sob
aprovação de comitê de ética, produzindo de propósito os espécimes de erro de que o bloco de detecção
precisa. O segundo é executar a comparação entre modelos já preparada no repositório, que testaria a
hipótese pelo lado oposto, pois se a verificação simbólica cumpre o que se afirma, um modelo mais fraco
e mais barato deveria errar mais e ainda assim convergir com mais iterações, e se não convergir terá
sido encontrado o piso de capacidade da arquitetura. Ambos os resultados interessam ao professor sem
orçamento, para quem um modelo aberto executado localmente seria o caminho natural. O terceiro é
atacar o ponto cego revelado pela Q061, por exemplo pedindo a um segundo agente que formalize o mesmo
enunciado de forma independente e comparando as duas formalizações antes da verificação. Há ainda
desdobramentos de alcance maior, como gerar figuras para as habilidades de conversão entre registros,
diversificar contextos e dados numéricos entre ciclos, ampliar o catálogo para as demais habilidades
do Ensino Médio que admitem verificação simbólica, adotar um \emph{framework} de orquestração caso o
fluxo passe a exigir paralelismo entre avaliadores e, mais adiante, investigar os efeitos do uso
dessas questões em sala de aula, o que esta pesquisa deliberadamente não fez.

Uma última consideração diz respeito ao uso responsável da inteligência artificial na Educação
Matemática. Os modelos de linguagem já estão nas mãos dos professores, e a questão relevante deixou de
ser se eles serão usados para passar a ser como. Este trabalho procurou mostrar que é possível
construir em torno desses modelos mecanismos que tornem os seus erros visíveis, que declarem com
honestidade o que foi e o que não foi conferido e que devolvam ao professor, com informação melhor, a
decisão que sempre foi dele. Uma ferramenta assim não diminui a responsabilidade docente, mas permite
que ela seja exercida onde de fato importa.
